% chapters/ch19-decomposition-heuristics.tex
\chapter{Decomposition Heuristics}

\begin{learningobjectives}
In this chapter you will learn how to split relations by meaning.
You will recognize mixed-grain structures and repeating groups.
You will practice decomposition without relying on convenience.
\end{learningobjectives}

\section{One relation, one meaning}
A relation should describe one kind of fact.
If it mixes meanings, it invites contradictions.
This is the core heuristic.

\begin{principlebox}
\textbf{Principle}\\
Split by meaning, not by convenience.
\end{principlebox}

\section{Repeating groups and collections}
If an attribute holds many values, the meaning is unclear.
Lists hide multiplicity.
They belong in a separate relation.

\subsection{Multi-valued attributes}
If an entity can have many phone numbers, the phone numbers are a collection.
Collections deserve their own identity and constraints.

\begin{examplebox}
\textbf{Repeating group}\\
\rel{Member}(member\_id, name, phone1, phone2, phone3)\\
Better meaning: \rel{Member}(member\_id, name) and \rel{MemberPhone}(member\_id, phone).
\end{examplebox}

\section{Mixed grain}
A mixed-grain relation blends entity facts with event facts.
This causes duplication and update anomalies.

\subsection{Entity versus event}
A person’s name is stable.
An appointment time is an event.
Do not store them in the same relation unless the event is the entity.

\begin{examplebox}
\textbf{Mixed grain}\par\smallskip
\[
\rel{Appointment}(appointment\_id, start\_time, doctor\_name, patient\_name)
\]
Better meaning:
\[
\begin{aligned}
\rel{Doctor}(&doctor\_id, name), \\
\rel{Patient}(&patient\_id, name), \\
\rel{Appointment}(&appointment\_id, start\_time, \\
&doctor\_id \rightarrow \rel{Doctor}, patient\_id \rightarrow \rel{Patient})
\end{aligned}
\]
\end{examplebox}

\section{Derived data}
Derived attributes can be useful, but they are not primary facts.
If you include them, be explicit that they are derived.
Otherwise they will be mistaken for ground truth.

\section{Stability and change}
If part of a relation changes at a different rate, you likely have two meanings.
Split slow-changing facts from fast-changing facts.
This keeps history consistent.

\subsection{Time-varying attributes}
When a value changes over time, consider a separate relation to track history.
This is a conceptual decision before any SQL.

\section{Decomposition checklist}
Ask three questions:
\begin{itemize}
  \item Does each attribute describe the same kind of fact?
  \item Can the attributes change together?
  \item Would a missing subset still make sense?
\end{itemize}
If any answer is no, decomposition is likely needed.

\begin{chaptersummary}
Decomposition is about meaning, not performance.
Repeating groups, mixed grain, and time-varying facts signal a need to split.
A single relation should represent a single kind of fact.
\end{chaptersummary}

\begin{exercisebox}
Choose a small domain and draft a single relation that mixes facts.
\begin{enumerate}[label=\arabic*.]
  \item Identify at least two different meanings in the relation.
  \item Decompose it into two or more relations with clear meaning.
  \item Describe the relationships between the new relations.
  \item State one constraint that is easier to express after decomposition.
  \item Explain how the decomposition improves clarity.
\end{enumerate}
\end{exercisebox}
